% !!! STALE — DO NOT COMPILE FOR SUBMISSION !!!
% Every result number in this file was computed from a benchmark pipeline
% confirmed to leak ground_truth_claims into error_logs
% (self_consistency.py:120, evaluation.py:334) and into the seeded
% Neo4j/Qdrant evidence itself (seeding.py), uniformly across all three
% context conditions. Do not build submission.zip from this file until
% the fix roadmap completes and every number here has been replaced from
% a re-run on corrected data.
\documentclass{article}

% Use numeric citations (matches the plain \bibitem entries below, which
% have no author-year optional label).
\PassOptionsToPackage{numbers}{natbib}

% Submission-time options: double-blind workshop review, no "final".
\usepackage[dblblindworkshop]{neurips_2026}
\workshoptitle{Who Verifies the Agents?}

\usepackage[utf8]{inputenc}
\usepackage[T1]{fontenc}
\usepackage{hyperref}
\usepackage{url}
\usepackage{booktabs}
\usepackage{amsfonts}
\usepackage{amsmath}
\usepackage{nicefrac}
\usepackage{microtype}
\usepackage{xcolor}
\usepackage{graphicx}

\title{Does More Agents Mean More Verification? A Matched-Compute Study of
Multi-Agent Hallucination and Graph-Grounded Claim Verification}

\author{%
  Anonymous Author(s) \\
  Anonymous Institution \\
  \texttt{anonymous@example.com} \\
}

\begin{document}

\maketitle

\begin{abstract}
Multi-agent LLM systems are increasingly deployed for real-world diagnostic
tasks, on the implicit assumption that decomposing a problem across
specialist agents improves both accuracy and trustworthiness. We test this
assumption on a real (non-simulated) multi-agent root-cause-analysis (RCA)
system built over a live Kubernetes knowledge graph, and introduce
Graph-Provenance Claim Scoring (GPCS), an evidence-grounded verification
mechanism that scores each atomic claim in a generated report by its
semantic and graph-topological proximity to retrieved evidence, contrasted
against a self-consistency baseline. Across 25 real incident scenarios we
find: (1) GPCS and self-consistency disagree on 36\% of claims, indicating
the two verification signals have materially different blind spots rather
than being redundant; (2) structured, ranked retrieval measurably
outperforms both no-context and unranked raw-context conditions; and (3),
most importantly for verifying multi-agent systems specifically, a
matched-compute control shows the real 5-specialist-agent architecture is
\emph{more} hallucination-prone than a single LLM sampled the same number
of times (44.2\% vs.\ 31.5\% unsupported-claim rate, paired 95\% CI
$[+0.053, +0.201]$, Wilcoxon $p=0.0018$) --- a statistically significant
negative result. We argue this demonstrates that verifying multi-agent
systems requires scrutinizing the aggregation step itself, not just the
individual agents, and that no single verification signal available to us
caught this failure mode in isolation.
\end{abstract}

\section{Introduction}
\label{sec:intro}

Multi-agent LLM architectures are being deployed for increasingly
consequential, real-world tasks --- code review, financial analysis,
infrastructure incident response --- under the assumption, often implicit
and rarely tested against a compute-matched control, that decomposing a
task across specialist agents with a consensus step improves both accuracy
and trustworthiness over a single model call. Separately, a growing body
of work asks how such systems should be \emph{verified}: current
evaluation practice largely tests isolated models rather than interactive,
multi-agent settings, and even the auditing mechanisms proposed for
agentic systems can themselves be corrupted or miscalibrated, since they
are frequently themselves LLM-based agents subject to the same failure
modes they are meant to detect~\citep{aigi2026,anthropic2026}.

This paper reports on a real, deployed multi-agent system --- which we
refer to as the \emph{RCA system} throughout to preserve anonymity --- a
GraphRAG-powered root-cause-analysis platform for Kubernetes incidents,
built and evaluated not as a simulation but against live infrastructure
telemetry, a real knowledge graph, and real LLM calls throughout. We use
it as a testbed for two questions directly relevant to verifying agentic
systems:

\textbf{RQ1 (verifying claims).} Given a generated diagnostic report, does
an evidence-grounded verification signal (GPCS) agree with a purely
behavioral one (self-consistency across resampled generations), or do
they catch different failures? If they diverge substantially, relying on
either one alone under-verifies the system.

\textbf{RQ2 (verifying architecture).} Does a real multi-agent
architecture (five specialist agents plus a consensus step) actually
reduce hallucination relative to a matched-compute single-LLM baseline, or
does the added structure introduce new, unverified failure surface ---
specifically, an unaudited aggregation step where a single weak
specialist's error can be woven into the final consensus without
correction?

We report both as measured, including where the results are unflattering
to the system we built: the real 5-agent architecture is significantly
worse than its own matched-compute single-LLM baseline on this benchmark.
We believe this negative result is itself the paper's main contribution to
the ``who verifies the agents'' question --- it is direct empirical
evidence that adding agents is not verification-neutral, and that the
aggregation step of a multi-agent system needs to be an explicit object of
verification, not an assumed-safe glue layer.

\textbf{Contributions.}
\begin{enumerate}
    \item GPCS, an evidence-grounded claim-verification mechanism for
    graph-backed RAG systems, contrasted head-to-head against a
    self-consistency baseline on the same 1{,}685 scored claims
    (Sections~\ref{sec:system},~\ref{sec:results-rq1}).
    \item A controlled, three-condition retrieval ablation (no context /
    unranked raw context / ranked GraphRAG context) isolating whether
    structured retrieval earns its complexity for claim grounding
    (Section~\ref{sec:results-retrieval}).
    \item A matched-compute control isolating whether a real multi-agent
    architecture's benefit (if any) comes from genuine multi-agent
    structure or is confounded with simply making more LLM calls --- with
    a negative result reported as measured
    (Section~\ref{sec:results-matched-compute}).
    \item All results are from a real, non-simulated system: real
    Kubernetes telemetry, a real Neo4j/Qdrant-backed knowledge graph, and
    real LLM API calls logged end-to-end, not a synthetic or
    offline-replay benchmark.
\end{enumerate}

\section{Related Work}
\label{sec:related}

\textbf{GraphRAG and structured retrieval.} GraphRAG~\citep{edge2024} and
follow-on work typically assume a static, once-built graph over a fixed
document corpus. Our setting differs in that the graph is continuously
mutating operational telemetry, and retrieval already incorporates a
temporal-window constraint and recency-decay term not present in the
canonical formulation. We do not centre this distinction in the present
paper --- it is out of scope for a verification-focused workshop --- but
note it as context for why our retrieval ablation
(Section~\ref{sec:results-retrieval}) is a live-system comparison, not a
replay over a fixed corpus.

\textbf{Hallucination detection.} The dominant self-consistency
family~\citep{wang2023}, sample repeatedly and flag claims that do not
recur, is model-internal and domain-agnostic. GPCS is evidence-grounded
rather than model-internal, and specific to settings --- like ours ---
where a structured evidence graph already exists as a byproduct of the
system rather than as an added verification cost. We are not aware of
prior work directly comparing an evidence-grounded and a
self-consistency-based verifier head-to-head on the same claim set with
paired statistical testing, which is the comparison
Section~\ref{sec:results-rq1} provides.

\textbf{Multi-agent LLM systems and matched-compute controls.} Surveys of
multi-agent LLM systems~\citep{guo2024} note that few papers claiming a
multi-agent benefit control for the ``more LLM calls = better'' confound.
Application papers claiming multi-agent benefits in operational domains in
particular rarely compare against a matched-compute single-agent or
ensemble baseline. Our matched-compute control
(Section~\ref{sec:results-matched-compute}) is a direct, real-system
instance of this missing comparison.

\textbf{Verifying agentic systems.} Current evaluation practice for
agentic systems is reported to concentrate on isolated-model benchmarks,
with interactive multi-agent evaluation ``underdeveloped''~\citep{aigi2026}.
Separately, recent work on agentic misalignment has surfaced a specific
failure mode directly relevant to this paper: an auditing agent can
itself produce a corrupted evaluation, e.g.\ by sharing the objection of
the agent it is meant to be auditing and declining to flag
it~\citep{anthropic2026}. Our matched-compute result is an instance of the
same underlying concern from a different angle --- not an auditor being
corrupted, but a \emph{consensus} step compounding rather than correcting
specialist error, which is invisible unless a matched-compute control is
explicitly run.

\section{System}
\label{sec:system}

The RCA system ingests Kubernetes telemetry (pod/deployment/node state,
logs, metrics, events) into a Neo4j property graph, indexes evidence
documents in Qdrant for dense retrieval, and answers incident-investigation
queries via three retrieval modes: keyword (lexical), vector (dense
semantic), and hybrid (vector similarity + graph hop-proximity + recency
decay, i.e.\ GraphRAG). Investigation requests are handled by five
LLM-backed specialist agents (monitoring, logs, deployment, topology,
security), each producing an independent finding, combined by a
static-weight consensus engine into a single report (title, cause,
recommendation, severity, evidence). This multi-agent architecture and the
claim-verification mechanisms described below are, together, the core
method of this paper: the LLM specialists, the LLM-based claim extractor,
and the self-consistency sampler are non-standard, original components of
the approach, not incidental implementation choices.

\subsection{Graph-Provenance Claim Scoring (GPCS)}
\label{sec:gpcs}

Given a generated report, GPCS: (A) extracts atomic factual claims via a
structured-output LLM call, typed as \texttt{temporal}, \texttt{causal},
\texttt{entity\_relationship}, or \texttt{state}; (B) for each claim,
retrieves candidate supporting evidence via semantic search over the claim
embedding and graph traversal seeded from any named entities in the
claim; (C) scores each claim as
\begin{multline}
    \text{trust\_score} = w_1 \cdot \text{semantic\_alignment}
    + w_2 \cdot \text{graph\_proximity} \\
    + w_3 \cdot \text{source\_reliability}
    - w_4 \cdot \text{path\_length\_penalty},
    \label{eq:gpcs}
\end{multline}
where \text{semantic\_alignment} is cosine similarity to the
best-matching evidence embedding, $\text{graph\_proximity} =
1/(1+\text{hop\_distance})$, \text{source\_reliability} is a fixed
per-evidence-type weight (metric/event readings weighted above free-text
logs), and \text{path\_length\_penalty} --- the one component with no
analogue in the system's existing retrieval-ranking formula --- discounts
claims resting on long inferential chains even when the final hop is a
strong match. Claims scoring below a calibrated threshold (semantic-floor
calibrated at 0.30 against live query score distributions, distinguishing
vague claims topping out at 0.16--0.30 from genuine matches at
0.33--0.87) are labeled unsupported. (D) Report-level
\texttt{unsupported\_claim\_rate} is aggregated overall and by claim type.

\subsection{Self-consistency baseline}
\label{sec:selfconsistency}

For the same incident, the system generates 3 samples at elevated
temperature, extracts claims from each identically to GPCS Step A, and
flags a claim from the primary generation as unsupported if it does not
recur (by semantic match) across the other samples. This is model-internal
and requires no evidence graph, making it the natural behavioral-signal
contrast to GPCS's evidence-grounded signal.

\section{Experimental Setup}
\label{sec:setup}

\textbf{Benchmark.} 25 real incident scenarios spanning 5 categories
(Kubernetes, networking, security, deployment, observability), each with
ground-truth expected evidence tags, seeded into a live cluster and graph.

\textbf{Conditions.} Each scenario is run under 3 retrieval-context
conditions --- \texttt{none} (agents reason from raw error logs alone, no
retrieval), \texttt{raw} (every seeded evidence item, concatenated,
unranked), and \texttt{hybrid} (GraphRAG's own ranked retrieval) --- giving
75 (scenario $\times$ condition) generation-and-scoring passes.

\textbf{Matched-compute control.} For the multi-agent-vs-single-LLM
comparison (Section~\ref{sec:results-matched-compute}), the \emph{same}
hybrid-retrieved evidence is given to two arms per scenario: the real
5-agent-plus-consensus system (6 LLM calls: 5 specialists + 1 consensus),
and a single LLM given the same evidence, sampled 5 times and
self-consistency-checked (5 LLM calls), with the primary sample's claims
scored by the identical GPCS instance used for the agent arm --- the
verification method is held constant, only generation architecture and
call count vary (150 vs.\ 125 total calls across 25 scenarios --- close,
not identical; we report this gap rather than claim exact parity).

\textbf{Statistics.} All headline deltas are evaluated with paired
bootstrap confidence intervals (10{,}000 resamples, seed=42) and the
Wilcoxon signed-rank test, computed per-scenario (or per
scenario$\times$condition pair) to respect the paired structure of the
design, given the small sample ($n=25$ scenarios; $n=75$ for the pooled
scenario$\times$condition comparison).

\textbf{Provenance.} Every LLM request/response for this run is logged
end-to-end (credentials redacted); this is a live-system run, not a replay
over cached model outputs.

\section{Results}
\label{sec:results}

\subsection{GPCS vs.\ self-consistency: divergent verification signals (RQ1)}
\label{sec:results-rq1}

Across all 75 scenario$\times$condition passes, 1{,}777 claims were
extracted; GPCS scored 1{,}685 of them (94.8\% --- the remainder is a
known claim-segmentation mismatch between two independent extraction
calls, not a correctness bug, discussed in Section~\ref{sec:limitations}).
Overall agreement between GPCS and self-consistency: \textbf{1{,}079/1{,}685
(64.0\%)}.

The two methods are not simply redundant. GPCS flags 43.7\% of claims as
unsupported; self-consistency flags 51.5\% --- self-consistency is the
stricter signal, with median recurrence rate of 0.0 (more than half of all
claims never recur verbatim across 3 samples at all). GPCS is more
lenient: any claim clearing the $\sim$0.30 semantic-relevance floor with
reasonable graph proximity passes, regardless of whether it happens to
recur across samples. Of the 606 cases where the two methods disagree,
60.9\% are GPCS-supported-but-self-consistency-unsupported, and 39.1\% the
reverse --- each method independently catches cases the other misses, in
both directions, not just one.

Paired bootstrap over all 75 groups: mean delta $-0.087$, 95\% CI
$[-0.119, -0.055]$ (excludes zero), Wilcoxon $p\approx0.0000$ --- the
divergence is statistically robust, not sampling noise.

By claim type, agreement clusters at 63--65\% for
causal/entity-relationship/state claims, reaches 76.9\% for temporal
claims ($n=26$, small sample), and is lowest, 60.4\%, for \texttt{general}
process-commentary claims (e.g.\ ``the finding is inconclusive''), which
are inherently hardest to ground since they concern the investigation
process itself rather than a checkable fact.

\begin{figure}[t]
    \centering
    \includegraphics[width=0.48\linewidth]{figures/unsupported_rate_by_claim_type.png}
    \hfill
    \includegraphics[width=0.48\linewidth]{figures/agreement_heatmap.png}
    \caption{Left: GPCS vs.\ self-consistency unsupported-claim rate by
    claim type. Right: agreement/disagreement cross-tabulation. Neither
    signal dominates the other across claim types.}
    \label{fig:claimtype-heatmap}
\end{figure}

\textbf{Implication for verification design.} Neither signal dominates. A
system relying on only one --- the common case in current practice ---
will silently miss a substantial fraction, comparable in size to what it
catches, of the failures the other signal would flag.

\subsection{Structured retrieval earns its complexity for claim grounding (RQ1, supporting)}
\label{sec:results-retrieval}

\begin{table}[t]
    \centering
    \caption{Retrieval-context ablation. \texttt{hybrid} wins on every
    column; \texttt{raw} is worst on every column, including agreement
    with self-consistency (SC) --- worse than no retrieved context at all.}
    \label{tab:context-ablation}
    \small
    \setlength{\tabcolsep}{4.5pt}
    \begin{tabular}{lrrrr}
        \toprule
        Context condition & $n$ claims & GPCS/SC agree & GPCS unsupp. & SC unsupp. \\
        \midrule
        \texttt{hybrid} (ranked GraphRAG) & 578 & \textbf{66.1\%} & \textbf{42.0\%} & \textbf{47.9\%} \\
        \texttt{none} (no context)        & 550 & 65.1\%          & 42.5\%          & 54.5\%          \\
        \texttt{raw} (unranked dump)      & 557 & 60.9\%          & 46.5\%          & 52.2\%          \\
        \bottomrule
    \end{tabular}
\end{table}

Table~\ref{tab:context-ablation} shows unranked evidence is not neutral,
it is actively harmful: more raw material gives the generator more
surface area to produce claims that sound evidence-backed without being
tied to the actually-relevant evidence.

The paired, per-scenario hybrid-vs-raw delta does not reach significance
at $n=25$ (mean $+0.050$, 95\% CI $[-0.017, +0.114]$, Wilcoxon $p=0.15$)
--- we report the direction as real and measured but the magnitude as not
yet statistically settled. A separate, purely retrieval-quality metric
(tag recall, not claim-level grounding) does reach significance: hybrid
vs.\ keyword tag-recall delta $+0.150$, 95\% CI $[+0.060, +0.240]$,
$p=0.0055$ (Figure~\ref{fig:recall}).

\begin{figure}[t]
    \centering
    \includegraphics[width=0.55\linewidth]{figures/retrieval_recall.png}
    \caption{Tag recall by retrieval method, with 95\% bootstrap CI error
    bars.}
    \label{fig:recall}
\end{figure}

\subsection{Matched-compute control: a negative result for multi-agent architecture (RQ2 --- main finding)}
\label{sec:results-matched-compute}

\begin{table}[t]
    \centering
    \caption{Matched-compute control: real 5-agent consensus system vs.\ a
    single LLM sampled the same evidence 5 times. The single-LLM baseline
    is \emph{less} hallucination-prone despite fewer LLM calls.}
    \label{tab:matched-compute}
    \small
    \setlength{\tabcolsep}{4.5pt}
    \begin{tabular}{lrrr}
        \toprule
        Arm & LLM calls (25 scenarios) & Mean unsupp. rate & Median \\
        \midrule
        Real 5-agent consensus                   & 150 ($6\times25$) & \textbf{44.2\%} & 44.4\% \\
        Single-LLM (5-sample self-consistency)    & 125 ($5\times25$) & \textbf{31.5\%} & 31.2\% \\
        \bottomrule
    \end{tabular}
\end{table}

The single-LLM baseline had a \emph{lower} hallucination rate in
\textbf{19 of 25 scenarios} (agents won in 6, no ties). Paired
per-scenario delta (agents $-$ single-LLM): mean $+0.127$, 95\% CI
$[+0.053, +0.201]$ (excludes zero), Wilcoxon $p=0.0018$ --- significant,
not sampling noise.

On this benchmark, the 5-specialist-agent architecture does not earn its
complexity over a matched-compute single-LLM baseline --- it is measurably
more hallucination-prone, not less. We report this as measured, per our
own guardrail against adjusting the evaluation harness to flatter the
system we built.

A plausible, not yet confirmed, mechanism: the five specialists report
independently and are combined by a static-weight consensus vote with no
cross-agent critique or correction step. A single low-confidence or
mistaken specialist finding can be woven into the final narrative without
being caught or revised. The single-LLM baseline has no such aggregation
step to introduce compounding error --- it reasons over the same evidence
once, directly. This is, we argue, exactly the kind of failure mode
``verifying the agents'' needs to target: not the individual specialist
calls, which are individually unremarkable, but the \emph{aggregation
step}, which is where the architecture's own untested assumption (that
combining specialists helps) breaks down.

We note two limitations on this comparison specifically (see
Section~\ref{sec:limitations} for the full list): call counts are close
but not identical (150 vs.\ 125, a small residual compute advantage for
the single-LLM arm that does not plausibly explain a 12.7-point gap), and
claim-count asymmetry between arms was not separately controlled for in
this pass.

\subsection{Neuro-symbolic retrieval ablation}
\label{sec:results-neurosymbolic}

Re-framing the three retrieval modes by symbolic/neural character
(keyword = symbolic, vector = neural, hybrid = neuro-symbolic),
retrieval-quality scores (does the method surface $\geq$50\% of a
scenario's expected evidence tags) are near-ceiling for all three (vector
25/25, keyword 24/25, hybrid 24/25) --- hybrid does not clearly beat
either half on this metric, a result we report plainly rather than adjust
for. The qualitative pattern is more informative than the aggregate:
keyword's misses concentrate on paraphrased or synonym-shifted concepts
(e.g.\ missing ``memory;killed'' when the evidence says ``OOMKilled''
verbatim); vector generalizes past these but still struggles with
numeric/status-code-style tags (e.g.\ \texttt{http\_429}); hybrid does not
clearly correct keyword's specific misses, plausibly because its
recency-decay term can suppress otherwise-relevant older evidence. With
$n=25$ and only one scenario failing the retrieval-quality bar outright,
this is a suggestive pattern for a larger-sample follow-up, not a settled
finding.

\section{Limitations}
\label{sec:limitations}

\begin{itemize}
    \item \textbf{$n=25$ scenarios.} Every interval above is
    correspondingly wide; treat point estimates (e.g.\ ``44.2\% vs.\
    31.5\%'') as approximate even where the direction is statistically
    significant.
    \item \textbf{Single LLM provider} across all conditions in this run;
    cross-provider replication is future work.
    \item \textbf{5.2\% of claims lack a GPCS score} due to independent
    claim segmentation between two separate extraction calls
    (self-consistency's and GPCS's) at nonzero sampling temperature ---
    excluded from the joined agreement count, not miscounted, but
    shrinking the effective sample.
    \item \textbf{Matched-compute call counts are close, not identical}
    (150 vs.\ 125); we report the gap rather than claim exact parity.
    \item \textbf{GPCS's own verification is intentionally narrow} --- it
    asks whether a claim is supported by the specific evidence the system
    retrieved, not whether it is true in any general sense. This is a
    deliberate scope choice (Section~\ref{sec:gpcs}) but is itself a
    limitation on how far ``verified by GPCS'' can be read as ``verified,
    full stop'' --- directly relevant to a workshop on verifying agents,
    since GPCS is itself an imperfect verifier with its own known blind
    spot (Section~\ref{sec:results-rq1}'s 36\% disagreement rate with an
    independent signal is direct evidence of this).
\end{itemize}

\section{Discussion: implications for ``who verifies the agents''}
\label{sec:discussion}

Two results in this paper speak directly to the workshop's central
question. First, two independent, reasonable verification signals for the
same claims disagree more than a third of the time --- verifying an
agentic system's outputs with only one signal available today leaves a
comparably-sized blind spot, whichever one is chosen. Second, and more
consequentially, a real multi-agent architecture that was assumed, by its
own designers, prior to this control, to be more trustworthy than a
single model turned out, under a fair matched-compute comparison, to be
significantly \emph{less} trustworthy --- and no per-agent inspection
would have surfaced this, since each specialist's individual output is
unremarkable; the failure is emergent in the aggregation step. We think
this is a concrete, empirical argument that verifying multi-agent systems
cannot stop at verifying the agents' individual outputs or even the final
report's claims --- it has to include a matched-compute, architecture-level
check of whether the multi-agent structure itself is earning its
complexity, run as a standard part of evaluating any such system before
deployment.

\textbf{Broader impacts.} A more reliable, evidence-grounded verification
layer for automated incident-response tooling has a direct positive
application: fewer unverified, plausible-sounding but ungrounded claims
reaching an on-call engineer who must act on them quickly. The main risk
this work highlights is the opposite failure: over-trusting an imperfect
verifier (Section~\ref{sec:limitations}) or an unaudited multi-agent
consensus step in a high-stakes production setting, which is precisely
the failure mode our matched-compute result demonstrates can happen
silently. We see this as an argument for routinely running matched-compute
and multi-signal verification controls before deploying agentic tooling
in consequential settings, not a reason to avoid such tooling altogether.

\section{Conclusion}
\label{sec:conclusion}

We presented a real, deployed multi-agent RCA system, an evidence-grounded
claim verifier (GPCS) contrasted against a self-consistency baseline, and
--- the paper's central empirical contribution --- a matched-compute
control showing that adding agents made this particular system measurably
more hallucination-prone, not less. We report this as a genuine negative
result, consistent with treating honest measurement as the actual
deliverable of this line of work rather than a system-promotion exercise.
We hope both the method (matched-compute controls as a standard
verification tool for multi-agent claims) and the finding (aggregation,
not just agent quality, needs verification) are useful to the broader
question this workshop is organized around.

\bibliographystyle{plain}
\begin{thebibliography}{9}

\bibitem{edge2024}
Darren Edge, Ha Trinh, Newman Cheng, Joshua Bradley, Alex Chao, Apurva
Mody, Steven Truitt, and Jonathan Larson.
\newblock From local to global: A graph {RAG} approach to query-focused
summarization.
\newblock \emph{arXiv preprint arXiv:2404.16130}, 2024.

\bibitem{wang2023}
Xuezhi Wang, Jason Wei, Dale Schuurmans, Quoc Le, Ed Chi, Sharan Narang,
Aakanksha Chowdhery, and Denny Zhou.
\newblock Self-consistency improves chain of thought reasoning in language
models.
\newblock \emph{International Conference on Learning Representations
(ICLR)}, 2023. arXiv:2203.11171.

\bibitem{guo2024}
Taicheng Guo, Xiuying Chen, Yaqi Wang, Ruidi Chang, Shichao Pei, Nitesh~V.
Chawla, Olaf Wiest, and Xiangliang Zhang.
\newblock Large language model based multi-agents: A survey of progress
and challenges.
\newblock \emph{International Joint Conference on Artificial Intelligence
(IJCAI)}, 2024. arXiv:2402.01680.

\bibitem{aigi2026}
Oxford Martin AI Governance Institute.
\newblock Open problems in frontier {AI} risk management.
\newblock 2026.

\bibitem{anthropic2026}
Anthropic Alignment Science.
\newblock Agentic misalignment in summer 2026.
\newblock \emph{Alignment Science Blog}, 2026.

\end{thebibliography}

\newpage
\section*{NeurIPS Paper Checklist}

\begin{enumerate}

\item {\bf Claims}
    \item[] Question: Do the main claims made in the abstract and introduction accurately reflect the paper's contributions and scope?
    \item[] Answer: \answerYes{}
    \item[] Justification: The abstract and Section~\ref{sec:intro} state exactly the three findings reported in Section~\ref{sec:results}, including the negative result, with the $n=25$ sample-size scope stated explicitly rather than implied.

\item {\bf Limitations}
    \item[] Question: Does the paper discuss the limitations of the work performed by the authors?
    \item[] Answer: \answerYes{}
    \item[] Justification: Section~\ref{sec:limitations} is a dedicated section covering sample size, single-provider scope, a known claim-segmentation mismatch, compute-count asymmetry in the matched-compute control, and the intentionally narrow scope of GPCS's own verification signal.

\item {\bf Theory assumptions and proofs}
    \item[] Question: For each theoretical result, does the paper provide the full set of assumptions and a complete (and correct) proof?
    \item[] Answer: \answerNA{}
    \item[] Justification: This is an empirical systems paper; it contains no theorems or formal proofs.

\item {\bf Experimental result reproducibility}
    \item[] Question: Does the paper fully disclose all the information needed to reproduce the main experimental results of the paper to the extent that it affects the main claims and/or conclusions of the paper (regardless of whether the code and data are provided or not)?
    \item[] Answer: \answerYes{}
    \item[] Justification: The GPCS scoring formula is given in full (Section~\ref{sec:gpcs}, Equation~\ref{eq:gpcs}), the self-consistency baseline is fully specified (Section~\ref{sec:selfconsistency}), and Section~\ref{sec:setup} specifies the benchmark, all three retrieval conditions, the matched-compute design and exact call counts, and the statistical methodology (paired bootstrap, 10{,}000 resamples, seed=42; Wilcoxon signed-rank).

\item {\bf Open access to data and code}
    \item[] Question: Does the paper provide open access to the data and code, with sufficient instructions to faithfully reproduce the main experimental results, as described in supplemental material?
    \item[] Answer: \answerNo{}
    \item[] Justification: The system's code and benchmark dataset are not linked at submission time to preserve double-blind anonymity. We intend to release both upon acceptance.

\item {\bf Experimental setting/details}
    \item[] Question: Does the paper specify all the training and test details (e.g., data splits, hyperparameters, how they were chosen, type of optimizer) necessary to understand the results?
    \item[] Answer: \answerYes{}
    \item[] Justification: Section~\ref{sec:setup} specifies the 25-scenario benchmark composition, all three context conditions, the matched-compute arms and call counts, and the GPCS semantic-relevance threshold and how it was calibrated (Section~\ref{sec:gpcs}).

\item {\bf Experiment statistical significance}
    \item[] Question: Does the paper report error bars suitably and correctly defined or other appropriate information about the statistical significance of the experiments?
    \item[] Answer: \answerYes{}
    \item[] Justification: All four headline deltas are reported with paired bootstrap 95\% confidence intervals (10{,}000 resamples, seed=42) and Wilcoxon signed-rank $p$-values (Sections~\ref{sec:results-rq1}--\ref{sec:results-matched-compute}), with the method stated in Section~\ref{sec:setup}.

\item {\bf Experiments compute resources}
    \item[] Question: For each experiment, does the paper provide sufficient information on the computer resources (type of compute workers, memory, time of execution) needed to reproduce the experiments?
    \item[] Answer: \answerYes{}
    \item[] Justification: All experiments call a hosted third-party LLM API rather than authors' own compute infrastructure, so the relevant resource metric is API call count, not GPU-hours; exact call counts are given for every comparison (150 vs.\ 125 for the matched-compute control, Section~\ref{sec:results-matched-compute}; 75 scenario$\times$condition passes for the main comparison, Section~\ref{sec:setup}).

\item {\bf Code of ethics}
    \item[] Question: Does the research conducted in the paper conform, in every respect, with the NeurIPS Code of Ethics \url{https://neurips.cc/public/EthicsGuidelines}?
    \item[] Answer: \answerYes{}
    \item[] Justification: The work uses telemetry from the authors' own test infrastructure, involves no human subjects, no crowdsourcing, and no personally identifiable or sensitive data.

\item {\bf Broader impacts}
    \item[] Question: Does the paper discuss both potential positive societal impacts and negative societal impacts of the work performed?
    \item[] Answer: \answerYes{}
    \item[] Justification: A dedicated Broader impacts paragraph closes Section~\ref{sec:discussion}, covering both the positive application (fewer ungrounded claims reaching on-call engineers) and the risk the paper's own negative result illustrates (over-trusting an imperfect verifier or unaudited multi-agent consensus in production).

\item {\bf Safeguards}
    \item[] Question: Does the paper describe safeguards that have been put in place for responsible release of data or models that have a high risk for misuse (e.g., pre-trained language models, image generators, or scraped datasets)?
    \item[] Answer: \answerNA{}
    \item[] Justification: The paper releases no model, generator, or scraped dataset with a high risk of misuse; it evaluates an internal diagnostic system against a benchmark of the authors' own seeded incident scenarios.

\item {\bf Licenses for existing assets}
    \item[] Question: Are the creators or original owners of assets (e.g., code, data, models), used in the paper, properly credited and are the license and terms of use explicitly mentioned and properly respected?
    \item[] Answer: \answerYes{}
    \item[] Justification: GraphRAG, self-consistency, and the multi-agent survey this work builds on and compares against are cited (Section~\ref{sec:related}); the embedding model and retrieval infrastructure used are standard, openly licensed open-source components.

\item {\bf New assets}
    \item[] Question: Are new assets introduced in the paper well documented and is the documentation provided alongside the assets?
    \item[] Answer: \answerNo{}
    \item[] Justification: The 25-scenario incident benchmark is a new asset, described in Section~\ref{sec:setup}, but is not yet packaged as a standalone, separately documented release at submission time; we intend to release it alongside the code upon acceptance.

\item {\bf Crowdsourcing and research with human subjects}
    \item[] Question: For crowdsourcing experiments and research with human subjects, does the paper include the full text of instructions given to participants and screenshots, if applicable, as well as details about compensation (if any)?
    \item[] Answer: \answerNA{}
    \item[] Justification: The paper involves no crowdsourcing and no human subjects.

\item {\bf Institutional review board (IRB) approvals or equivalent for research with human subjects}
    \item[] Question: Does the paper describe potential risks incurred by study participants, whether such risks were disclosed to the subjects, and whether Institutional Review Board (IRB) approvals (or an equivalent approval/review based on the requirements of your country or institution) were obtained?
    \item[] Answer: \answerNA{}
    \item[] Justification: The paper involves no human subjects research.

\item {\bf Declaration of LLM usage}
    \item[] Question: Does the paper describe the usage of LLMs if it is an important, original, or non-standard component of the core methods in this research? Note that if the LLM is used only for writing, editing, or formatting purposes and does \emph{not} impact the core methodology, scientific rigor, or originality of the research, declaration is not required.
    \item[] Answer: \answerYes{}
    \item[] Justification: LLMs are a core, non-standard component of the method itself, not a writing aid: the five specialist agents, the consensus step, GPCS's claim extraction, and the self-consistency baseline's sampling are all LLM-backed and fully described in Section~\ref{sec:system}.

\end{enumerate}


\end{document}
